% Source: https://artofproblemsolving.com/community/c3h2489361_some_fun_problems
% Stable identity: aops-topic-2489361-q02
Let $n=10a+b$ with digits $a\ge1$ and $b$. The condition $ab=\dfrac{10a+b}{2}$ becomes $2ab=10a+b$, i.e. $b(2a-1)=10a$. Since $\gcd(2a-1,a)=1$ (a common divisor divides $2a-(2a-1)=1$), $2a-1$ must divide $10$; as $2a-1$ is odd, $2a-1\in\{1,5\}$.
\begin{itemize}
\item $2a-1=1$: $a=1$, $b=10$, not a digit.
\item $2a-1=5$: $a=3$, $b=6$.
\end{itemize}
So $n=36$ (check: $3\cdot6=18=\dfrac{36}{2}$), and the digit sum is $\boxed{9}$.
